For a finite set $A\subset\mathbb{R}^2$, define
\[
\delta(A):=\min\{\|x-y\|:x,y\in A,\ x\ne y\}\quad\text{(minimum pairwise distance)},
\]
\[
\operatorname{diam}(A):=\max\{\|x-y\|:x,y\in A\}\quad\text{(diameter)}.
\]
For each $n\ge 2$, let
\[
D(n):=\inf\{\operatorname{diam}(A): A\subset\mathbb{R}^2,\ |A|=n,\ \delta(A)\ge 1\}.
\]
By scaling, this is equivalent to restricting to $\delta(A)=1$.
A set $A$ is \emph{diameter-minimising} if $|A|=n$, $\delta(A)\ge 1$, and $\operatorname{diam}(A)=D(n)$.

\medskip
\noindent\textbf{Question (\#99).} For all sufficiently large $n$, must every diameter-minimising $A$ contain three points that form an equilateral triangle of side length $1$?

\item
